\chapter{第二十章：人机协作技能——与AI共事的新基本功}

# 第二十章：人机协作技能——与AI共事的新基本功

\textit{"工具会改变我们使用它的方式，正如我们使用工具，工具也在使用我们。"}
——马歇尔·麦克卢汉

\hline\newpage
---

老王我老婆有个毛病，就是太依赖导航了。有一次我们开车去一个地方，她开着导航，结果导航指了一条路，导航说"前方三公里后左转"，结果那条路在修路，完全封死了。她愣在那里，不知道怎么办。我说"那就绕路啊"，她说"可是导航没说要绕路啊"。

我当时就想说一句：你到底是导航的主人，还是导航的奴隶？

后来我反思了一下自己，我发现我也有这毛病——只是不是导航，是搜索引擎。以前没有搜索引擎的时候，我脑子里记了很多东西：电话、地址、人名、典故。现在有了搜索引擎，我什么都不记了，因为"搜一下就知道了"。但问题是，搜一下不等于你知道，你只是知道"搜到过"。

AI来了，我估计我老婆的新毛病就是太依赖AI：让她写个报告，她说"让AI写"；让她分析个数，她说"让AI分析"；让她做个决策，她说"让AI建议"。

这不叫"用AI"，这叫"被AI用"。

\section{一、一个令人担忧的发现}

## 一、一个令人担忧的发现

在2024年的一项研究中，斯坦福大学的人机交互研究员发现了一个令人担忧的现象：

当人类长期依赖AI系统完成某些任务后，一旦AI系统不可用，人类的表现往往比那些从未使用过AI的对照组更差。

这个现象被称为"自动化依赖效应"。

更令人担忧的是，这个效应在AI助手（Copilot、ChatGPT等）用户中也开始出现。当人们习惯了让AI帮助写作、编程、分析后，一旦没有AI可用，他们发现自己无法完成同样的任务——即使这些任务是他们的"本职工作"。

这个发现有一个重要的含义：\textbf{与AI协作的技能，需要被刻意培养。否则，人类可能成为AI的依赖者，而不是AI的主导者。}

就像你天天吃外卖，突然有一天外卖停了，你发现你自己连煮个面都不会。你不是吃饭的主人，你是外卖平台的奴隶。

\section{二、为什么人机协作需要技能}

## 二、为什么人机协作需要技能

我们通常假设，使用AI工具是不需要技能的——你只要会打字，就能用ChatGPT。

但实际情况远比这复杂。

\textbf{用同样的AI工具，不同的人产生的结果质量差异巨大。}

有些人让AI写出的报告清晰，专业、有洞见。有些人让AI写出的报告平淡、泛泛、缺乏深度。差别不在于AI工具本身，而在于使用AI的人的技能。

这些技能包括：

\begin{itemize}
\item \textbf{任务定义}：知道什么任务适合交给AI，什么任务应该自己做
\item \textbf{提示工程}：知道如何给AI清晰的指令，如何引导AI产生更好的输出
\item \textbf{输出评估}：知道如何判断AI的输出是否正确、合理、有价值
\item \textbf{迭代优化}：知道如何基于AI的输出进行修改和改进

\end{itemize}
这些技能不是天生的。它们需要被学习和练习。

就像你会用筷子不等于你会中餐料理——筷子只是工具，料理是技能。你可以用筷子夹起任何东西，但只有技能才能让你做出一桌好菜。

\section{三、人机协作技能的四根支柱}

## 三、人机协作技能的四根支柱

人机协作技能，可以分为四根支柱。

\textbf{支柱一：任务定义}

任务定义是决定什么任务应该交给AI，什么任务应该保留给人类的能力。

这个能力看似简单，实际上需要深度的判断。

适合交给AI的任务通常是：信息检索、内容生成、数据分析、模式识别、标准化流程。

适合保留给人类的任务通常是：复杂问题的定义、价值判断、创新突破、人际关系、伦理决策。

但这个划分不是固定的——随着AI能力的进化，适合交给AI的任务范围在持续扩展。

有效的任务定义，需要持续追踪AI能力的进化，持续重新评估人机分工的边界。

\textbf{支柱二：提示工程}

提示工程（Prompt Engineering）是让AI产生更好输出的技术。

它的核心是：知道你想要什么，然后用清晰、准确的语言告诉AI。

但提示工程不只是"会说清楚"。它还包括：

\begin{itemize}
\item \textbf{提供足够的上下文}：让AI理解任务的背景和目标
\item \textbf{指定输出格式}：让AI知道你期望的输出形式
\item \textbf{引导AI的思考方向}：通过提示让AI朝正确的方向思考
\item \textbf{处理AI的局限}：知道AI的常见缺陷，在提示中预防它们

\end{itemize}
提示工程是一项可以被练习和改进的技能。随着你使用的AI工具越来越多，你可以积累更多的提示模板和最佳实践。

就像你请人帮忙，你得说清楚你要什么、什么时候要、有什么特别要求。你不能说"你懂的"，然后指望对方真的懂。

\textbf{支柱三：输出评估}

AI会犯错。AI会"幻觉"——生成听起来正确但实际错误的内容。AI的输出需要被评估。

输出评估的技能包括：

\begin{itemize}
\item \textbf{事实核查}：对AI输出的关键信息进行验证
\item \textbf{逻辑检验}：检查AI的推理过程是否有漏洞
\item \textbf{质量判断}：判断AI的输出是否有价值，是否满足任务需求
\item \textbf{偏差识别}：识别AI输出中可能存在的偏见或不当内容

\end{itemize}
很多人在使用AI时，跳过了输出评估这个环节。这是一种危险的倾向。AI的输出必须被审视，不能被盲目信任。

就像你网购，收了快递得检查一下东西对不对、坏没坏。你不能因为卖家说"保证质量"就什么都不看，到时候发现东西坏了再去找售后，耽误的是你自己的时间。

\textbf{支柱四：迭代优化}

AI很少能一次就产生完美的输出。好的结果，通常需要多轮的迭代优化。

迭代优化的技能包括：

\begin{itemize}
\item \textbf{解读AI的输出}：理解AI给出的内容，找到可以改进的方向
\item \textbf{给出改进指令}：用清晰的指令引导AI朝更好的方向改进
\item \textbf{整合AI的输出}：把AI的输出与自己的思考、其他的输入整合起来

\section{四、人机协作的常见陷阱}


\end{itemize}
在与AI协作时，有几个常见的陷阱需要避免。

\textbf{陷阱一：过度依赖}

过度依赖AI，丧失独立完成任务的能力。

这通常发生在那些把AI当作捷径而不是工具的人身上。他们用AI完成了任务，但不理解任务是如何被完成的。当AI不可用时，他们发现自己无法完成任务。

避免这个陷阱的方法是：在使用AI的同时，保持自己对任务的理解和控制。即使AI在做主要的工作，你也要理解它在做什么、为什么这样做。

\textbf{陷阱二：盲目信任}

盲目接受AI的输出，不进行批判性评估。

这通常发生在那些对AI过于信任或者懒得验证的人身上。AI会犯错，AI会幻觉，AI的输出可能包含偏见。盲目信任AI的输出，可能导致严重的错误。

避免这个陷阱的方法是：对AI的输出保持批判性审视，特别是那些涉及重要决策或敏感信息的内容。

\textbf{陷阱三：过于挑剔}

与过度依赖相对的另一个极端，是过于挑剔AI的输出，不断要求AI修改。

这可能发生在那些完美主义者身上。他们对AI的每一个输出都不满意，花大量时间在AI和自己之间来回修改，反而降低了效率。

避免这个陷阱的方法是：接受"足够好"的标准。AI的输出不是必须被修改到完美才能使用的。

\section{五、建立人机协作的实践}

## 五、建立人机协作的实践

培养人机协作技能，需要实践。

\textbf{实践一：设定人机协作的工作流程}

在你开始一个需要使用AI的工作之前，先想清楚：你打算如何使用AI？你打算把什么任务交给AI？你打算如何评估AI的输出？

把这种人机协作的工作流程显式化，可以帮助你更系统地思考人机协作。

\textbf{实践二：记录你的提示词和结果}

当你找到一个有效的提示词时，把它记录下来。当你发现某个提示词效果不好时，也记录下来。随着时间的推移，你会积累一个属于你自己的提示词库。

\textbf{实践三：定期反思AI的使用效果}

每隔一段时间，回顾一下你使用AI的工作：哪些地方AI真的帮到了你？哪些地方AI反而降低了效率？你对AI的使用是否有什么不当之处？

这种定期反思，可以帮助你持续改进你的人机协作技能。

\section{六、结语：人机协作是AI时代的新基本功}

## 六、结语：人机协作是AI时代的新基本功

人机协作技能，在不远的将来，会像今天的读写能力一样基础。

不会读写的人，在现代社会中几乎无法生存。不会与AI协作的人，在AI时代也会面临类似的问题。

但这不是一件可怕的事。这是人类适应新工具的又一次进化。

历史上，每一次重大技术革命都带来了新的"基本功"：工业革命带来了操作机器的能力，信息革命带来了使用计算机的能力。AI革命带来的，是人机协作的能力。

这种能力是可以被培养的。它需要刻意练习，需要反思，需要迭代。

那些能够学会与AI有效协作的人，将在AI时代占据优势。

麦克卢汉说"工具会改变我们使用它的方式，正如我们使用工具，工具也在使用我们"——这话在AI时代变得更加深刻。AI不只是一个工具，它是一个会反向塑造你的存在。你用AI的方式，决定了你是AI的主人还是AI的奴隶。

\hline\newpage
---

\textbf{本章小结：}

1. 自动化依赖效应：长期依赖AI后，一旦AI不可用，人类表现反而比从未使用AI的人更差
2. 人机协作需要技能：用同样的AI工具，不同人产生的结果质量差异巨大
3. 人机协作技能的四根支柱：任务定义（决定什么交给AI）、提示工程（让AI产生好输出）、输出评估（判断AI输出是否正确）、迭代优化（多轮改进结果）
4. 人机协作的三个陷阱：过度依赖（丧失独立能力）、盲目信任（忽略AI的错误）、过于挑剔（效率降低）
5. 建立人机协作实践：设定工作流程、记录提示词和结果、定期反思AI使用效果
6. 人机协作是AI时代的新基本功——像读写能力一样基础，可以被培养和迭代
